% !TeX program = xelatex
% 使用 XeLaTeX；Fandol 字体随 TeX Live 提供，无需系统字体。
\documentclass[UTF8,fontset=fandol,a4paper,zihao=5]{ctexart}
\usepackage{uestcdesign}
\linespread{1.12}
% 修改以下字段即可开始写作；含逗号的值请放在花括号内。
\uestcsetup{
theme=homework,
 type=课程作业,english={COURSE ASSIGNMENT / 01},
 title={线性代数 · 第一周作业},subtitle={矩阵、特征值与二次型},
 author=张同学,student-id=202600000000,
 school=信息与通信工程学院,course={高等线性代数},
 teacher=教师姓名,date={2026 年 10 月 13 日},short-title={第一周作业}
}
\begin{document}
\uestcheading
\begin{summarybox}[作业要求]
写出必要推导过程，给出结果并检查。此模板无独立封面，适合短作业、习题解答与随堂报告。
\end{summarybox}
\section{特征值与特征向量}
\textbf{题目：}求矩阵 $A$ 的特征值及对应特征向量。
\[
A=\begin{pmatrix}2&1\\1&2\end{pmatrix}.
\]
\textbf{解：}由特征多项式
\[
\det(A-\lambda I)=(2-\lambda)^2-1=(\lambda-1)(\lambda-3),
\]
得到 $\lambda_1=1$、$\lambda_2=3$。可分别取对应的单位特征向量
\[
q_1=\frac{1}{\sqrt{2}}\begin{pmatrix}1\\-1\end{pmatrix},\qquad
q_2=\frac{1}{\sqrt{2}}\begin{pmatrix}1\\1\end{pmatrix}.
\]
代入 $Aq_i=\lambda_iq_i$ 即可验证，且 $q_1^{\mathsf T}q_2=0$。
\section{正定性判别}
\textbf{题目：}判断上题矩阵是否正定，并说明理由。

\textbf{解：}矩阵 $A$ 为实对称矩阵，且两个特征值均大于零，所以 $A$ 正定。也可以通过顺序主子式验证：
\[
\Delta_1=2>0,\qquad\Delta_2=\det A=3>0.
\]
\section{学习记录}
\begin{itemize}
\item 实对称矩阵不同特征值对应的特征向量正交。
\item 特征向量归一化后，可以组成正交矩阵进行对角化。
\item 检查答案时，同时验证代数运算与概念使用条件。
\end{itemize}
\end{document}
